\section{Research taste y el \emph{Sutra del Diamante}: cómo juzgar si un problema vale la pena}

A la mitad de la entrevista hay un tramo largo dedicado a la research taste, en el que se menciona el verso del \emph{Sutra del Diamante} (《金刚经》) «como un sueño, una ilusión, una burbuja, una sombra». Esta parte no es un paréntesis místico, sino metodología de investigación. Lo que le interesa a Xie Saining (谢赛宁) es cómo mantener el juicio sobre el problema mismo en un entorno de investigación de gran incertidumbre, con una evaluación muy ruidosa y resultados de corto plazo que con frecuencia engañan.

\begin{figure}[H]
\centering
\includegraphics[width=0.92\textwidth]{research-taste-map.png}
\caption{Los componentes de la research taste: elección del problema, criterio estético, evidencia y momento oportuno.}
\end{figure}

\begin{knowledgebox}{Lectura de la figura: la research taste no es inspiración, sino un sistema de juicio a largo plazo}
Las cuatro esquinas de la figura son la elección del problema, el criterio estético, la evidencia y el momento oportuno. Para juzgar qué problema vale la pena, un investigador no puede fijarse solo en lo que está de moda ni en si puede publicar un artículo rápido; también debe preguntarse si el problema seguirá teniendo sentido dentro de diez años, si la evidencia respalda la dirección, si los recursos de la organización son suficientes y si él mismo puede tolerar una incertidumbre prolongada.
\end{knowledgebox}

\subsection{Criterio estético: contra las citas baratas y los eslóganes}

El párrafo anterior descompuso la research taste en problema, evidencia, criterio estético y momento oportuno; aquí se desarrolla primero el «criterio estético». No se trata de una preferencia de estilo, sino de la sensibilidad del investigador ante el abuso de conceptos, el empaque con eslóganes y la sustitución encubierta de la evidencia.

Al final, Xie Saining critica que se tome la frase de Wittgenstein (维特根斯坦) «los límites de mi lenguaje son los límites de mi mundo» para avalar directamente a los LLM, y también que se use a la ligera la frase de Feynman (费曼) «What I cannot create, I do not understand» para avalar los unified model. No se opone a la filosofía ni a las citas célebres, sino a las citas decorativas sacadas de contexto.

Detrás de esto está el criterio estético de la research taste: un argumento no puede sostenerse en citas de personajes célebres, sino en definiciones, supuestos, mecanismos y evidencia. La filosofía puede inspirar la investigación, pero no sustituye a un technical argument.

\begin{warningbox}{Una cita célebre no es un argumento técnico}
Poner una frase célebre de filosofía o de física al inicio de un artículo no demuestra por sí solo que la línea de modelos elegida sea correcta. Una investigación de calidad necesita explicar qué significa el concepto en su contexto original, si los supuestos siguen siendo válidos al trasladarlo a la IA y qué mecanismos y evidencia respaldan ese traslado.
\end{warningbox}

\subsection{Visión de largo plazo: del rechazo a la prueba del tiempo}

En la entrevista también se cuenta que un artículo fue rechazado en su momento por cuestiones de detalle, y que después se publicó en otro congreso y obtuvo un test-of-time award. La historia muestra que el valor de una investigación y el resultado de la revisión no siempre van de la mano. La revisión de corto plazo puede fijarse en exceso en el formato, los detalles o las preferencias dominantes del momento; el impacto de largo plazo depende más de que el problema sea real, de que el método sea reutilizable y de que la idea se incorpore a trabajos posteriores.

\begin{importantbox}{La escala de tiempo de la evaluación de la investigación}
La evaluación de corto plazo mira el accept / reject; la de largo plazo mira si la idea cambia la definición del problema, las herramientas metodológicas o el lenguaje de la comunidad. La verdadera research taste consiste en mantener, en medio del ruido de corto plazo, una dirección reutilizable a largo plazo.
\end{importantbox}

\subsection{Resumen de la sección}

La research taste es lo que une la línea técnica y la línea personal de este episodio. Explica por qué Xie Saining eligió vision, eligió FAIR, dejó un camino seguro y fundó una empresa para trabajar en world model, y explica también por qué le disgusta dejarse llevar por eslóganes y palabras de moda.
